\subsection{Interpretation of the Detectability Curve}

The central result of this work is the sigmoidal detectability curve
in Figure~\ref{fig:detectability_curve}. Above threshold,
$\beta \gtrsim 0.4$, the classifier reliably distinguishes GR from
amplitude-modulated waveforms with $\geq 95\%$ accuracy, even after
amplitude normalization. Feature-importance analysis shows that the
classifier uses tail percentiles of the amplitude distribution, which
is consistent with amplitude modulation concentrating energy near
merger. At threshold, $\beta \approx 0.25$, detection occurs with
probability approximately $0.75$ and ROC-AUC approximately $0.87$; the
tail statistics differ between classes by only a few percent,
approaching the intrinsic noise floor. Below threshold,
$\beta \lesssim 0.2$, detection fails completely.

\subsection{Relationship Between $\beta$ and Physical Beyond-GR Parameters}

The coefficient $\beta$ in Eqs.~\eqref{eq:amp_dev}--\eqref{eq:freq_dev}
is the amplitude of a specific quadratic-in-time modulation, not a
standard parametrized-GR coefficient. Direct comparison of our
threshold to physical constraints therefore requires care. In the
standard parametrized-test framework \citep{Yunes2019}, deviations are
expressed as fractional shifts $\delta \hat{\alpha}_k$ of post-Newtonian
phase coefficients, and current gravitational-wave observations bound
these shifts to roughly $|\delta \hat{\alpha}_k| \lesssim 10^{-3}$--$10^{-2}$
\citep{Abbott2021tests}. A controlled quadratic-in-time modulation of
the type studied here does not map onto a single PN coefficient, and
its effective $\beta$ differs from these bounds by a
waveform-dependent factor of order unity. Our threshold
$\beta_{\rm threshold} \approx 0.25$ should be read as a property of
the quadratic modulation adopted here; extending the analysis to the
standard parametrized-GR framework is a natural next step.

\subsection{Why the Negative Result Matters}

The failure at small $\beta$ is scientifically meaningful. It
quantifies the limits of ML-only detection. A common assumption in
ML-for-physics research is that neural networks will detect any signal
that is present. Our results show this is not the case for controlled
deviations in realistic noise. It also distinguishes ML detection
from matched filtering: matched-filter searches achieve detection
thresholds corresponding to $\rho_{\rm SNR} \sim 5$ by coherently
integrating the signal across its frequency evolution, whereas our
method operates on the whitened time-domain waveform directly. This
gap in sensitivity motivates hybrid approaches in which ML is
combined with physically motivated filtering.

\subsection{Comparison with Prior Work}

Prior ML work on gravitational waves has largely focused on detection
of GR signals in noise \citep{George2018, Gebhard2019} or on parameter
estimation \citep{Gabbard2020, Dax2021}. Beyond-GR searches with ML
have focused on simulated signals at high SNR or on clean data. Our
work extends this to a systematic study with real detector noise and
a quantified threshold.

\subsection{Limitations}

Several limitations should be noted. Our deviations follow a specific
quadratic-in-time form, $\propto \tau(t)^2$; other parametrizations
may have different detectability thresholds. The GW150914 study uses
a single template; generalization across event parameters has been
tested only in the injected-waveform studies. We have not tested
neutron-star or mixed binaries, where tidal effects provide additional
deviation channels. Finally, our pipeline does not include the
matched-filter step used in LIGO/Virgo analyses; adding it would
substantially improve sensitivity but would blur the boundary between
ML and classical detection methods.

\subsection{Future Work}

Several directions follow naturally. The classifier could be combined
with matched-filter outputs by using the matched-filter SNR time
series as input. The method could be applied to a population of
simulated binary mergers spanning the full LIGO parameter space.
Frequency-domain features such as spectrograms could be investigated
as alternatives to the time-domain approach. Probabilistic models
such as normalizing flows could be used to quantify classifier
uncertainty.
